% ============================================
% TABLA 1
% TABLA CORTA O CON POCA INFORMACIÓN
% ============================================

\begin{table}[htbp]
    \centering

    % Número de tabla
    \refstepcounter{table}\label{tab:intervalo-edades}
    \textbf{Tabla \thetable}

    \vspace{0.5\baselineskip}

    % Título de la tabla
    \textit{Ejemplo de tabla corta o con poca información}

    \vspace{0.5\baselineskip}

    \begin{tabular}{lccc}
        \toprule

        & \multicolumn{2}{c}{Género} & \\

        \cmidrule(lr){2-3}

        Intervalo de edades & F & M & Frecuencia \\

        \midrule

        1 a 2 años & 23 & 19 & 3 años \\

        3 a 4 años & 31 & 27 & 2 años \\

        5 a 6 años & 38 & 33 & 1 año \\

        7 a 8 años & 51 & 48 & 1 año \\

        \bottomrule
    \end{tabular}

    \vspace{0.5\baselineskip}

    \begin{flushleft}
        \textit{Nota.}
        Esta tabla muestra el intervalo de años y géneros
        encontrados en el estudio de caso realizado por el
        proyecto de estructura.

        \textit{Fuente.} Autor.
    \end{flushleft}

\end{table}